\documentclass[11pt]{article}

\usepackage[margin=1in]{geometry}
\usepackage[T1]{fontenc}
\usepackage{lmodern}
\usepackage{microtype}
\usepackage{setspace}
\usepackage{amsmath,amssymb}
\usepackage{hyperref}

\title{\textbf{Quantitative Retail: Mission, Method, and the Case for Evidence-Based Retail Trading}}\author{Michael Moody}\date{\today}

\begin{document}
\maketitle

\begin{abstract}
This paper sets out the foundational mission and methodological commitments of Quantitative Retail. Its purpose is to explain why the project exists, what standards it applies to retail trading claims, and why passive benchmarks must remain central to any serious evaluation of strategy performance.
\end{abstract}

\section{Introduction}
Quantitative Retail exists to test retail trading claims with data, explicit rules, realistic cost assumptions, and benchmark-aware analysis.

\section{Core Thesis}
Most publicly marketed retail trading strategies are not robust enough to justify strong confidence, and many retail traders would likely be better served by passive investing unless they can demonstrate a repeatable edge under realistic conditions.

\section{Method}
Claims will be evaluated through exact rule definition, out-of-sample testing, realistic implementation, regime analysis, and comparison against passive benchmarks.

\section{Conclusion}
Retail traders deserve evidence, not mythology.

\end{document}